\section{Формула Тейлора с остаточным членом в форме Лагранжа. Второе достаточное условие экстремума.}

\subsection{Формула Тейлора с остаточным членом в форме Лагранжа}

\begin{theorem}[Формула Тейлора с остаточным членом в форме Лагранжа]
    Пусть функция \(f(x)\) \(n+1\) раз дифференцируема на интервале \((a, b)\) и \(x_0, x \in (a, b)\). 
    Тогда \(\exists c \in (x_0; x)\) (или \(c \in (x; x_0)\)) такое, что
    \[
    f(x) = T_n(x) + \frac{f^{(n+1)}(c)}{(n+1)!} (x - x_0)^{n+1},
    \]
    где \(T_n(x) = \sum_{k=0}^n \frac{f^{(k)}(x_0)}{k!} (x - x_0)^k\) — многочлен Тейлора.
\end{theorem}

\begin{Proof}
    Зафиксируем \(x\) и рассмотрим вспомогательную функцию от переменной \(t\):
    \[
    \gamma(t) = f(x) - \sum_{k=0}^n \frac{f^{(k)}(t)}{k!} (x - t)^k - \frac{(x - t)^{n+1}}{(x - x_0)^{n+1}} R_n(x),
    \]
    где \(R_n(x) = f(x) - T_n(x)\) — остаточный член.
    
    Заметим, что:
    \begin{itemize}
        \item \(\gamma(x_0) = f(x) - T_n(x) - R_n(x) = f(x) - T_n(x) - (f(x) - T_n(x)) = 0\)
        \item \(\gamma(x) = f(x) - \sum_{k=0}^n \frac{f^{(k)}(x)}{k!} (x - x)^k - 0 = f(x) - f(x) = 0\)
    \end{itemize}
    
    Функция \(\gamma(t)\) непрерывна на \([x_0, x]\) и дифференцируема на \((x_0, x)\). 
    По теореме Ролля \(\exists c \in (x_0, x): \gamma'(c) = 0\).
    
    Вычислим производную \(\gamma'(t)\):
    \[
    \gamma'(t) = -\sum_{k=0}^n \left( \frac{f^{(k+1)}(t)}{k!} (x - t)^k - \frac{f^{(k)}(t)}{(k-1)!} (x - t)^{k-1} \right) + \frac{(n+1)(x - t)^n}{(x - x_0)^{n+1}} R_n(x).
    \]
    
    При упрощении все члены суммы, кроме последнего, сокращаются, и остаётся:
    \[
    \gamma'(t) = -\frac{f^{(n+1)}(t)}{n!} (x - t)^n + \frac{(n+1)(x - t)^n}{(x - x_0)^{n+1}} R_n(x).
    \]
    
    Подставляя \(t = c\) и используя \(\gamma'(c) = 0\), получаем:
    \[
    -\frac{f^{(n+1)}(c)}{n!} (x - c)^n + \frac{(n+1)(x - c)^n}{(x - x_0)^{n+1}} R_n(x) = 0.
    \]
    
    Сокращая на \((x - c)^n \neq 0\) (так как \(c \neq x\)), находим:
    \[
    R_n(x) = \frac{f^{(n+1)}(c)}{(n+1)!} (x - x_0)^{n+1}.
    \]
    
    Таким образом,
    \[
    f(x) = T_n(x) + \frac{f^{(n+1)}(c)}{(n+1)!} (x - x_0)^{n+1}.
    \]
\end{Proof}

\subsection{Второе достаточное условие экстремума}

\begin{theorem}[Второе достаточное условие экстремума]
    Пусть функция \(f(x)\) \(n\) раз дифференцируема в точке \(x_0\) и выполнено:
    \begin{enumerate}
        \item \(f'(x_0) = f''(x_0) = \dots = f^{(n-1)}(x_0) = 0\)
        \item \(f^{(n)}(x_0) \neq 0\)
    \end{enumerate}
    
    Тогда:
    \begin{itemize}
        \item Если \(n\) — чётное и \(f^{(n)}(x_0) > 0\), то \(x_0\) — точка локального минимума.
        \item Если \(n\) — чётное и \(f^{(n)}(x_0) < 0\), то \(x_0\) — точка локального максимума.
        \item Если \(n\) — нечётное, то \(x_0\) — точка перегиба (не экстремум).
    \end{itemize}
\end{theorem}

\begin{Proof}
    По формуле Тейлора с остаточным членом в форме Пеано (так как \(f^{(n)}(x_0)\) существует):
    \[
    f(x) = f(x_0) + \frac{f^{(n)}(x_0)}{n!} (x - x_0)^n + o((x - x_0)^n).
    \]
    
    Отсюда:
    \[
    f(x) - f(x_0) = \left( \frac{f^{(n)}(x_0)}{n!} + o(1) \right) (x - x_0)^n.
    \]
    
    Рассмотрим два случая:
    
    \textbf{Случай 1: \(n\) чётное.}
    Тогда \((x - x_0)^n > 0\) для всех \(x \neq x_0\). Знак разности \(f(x) - f(x_0)\) определяется знаком \(\frac{f^{(n)}(x_0)}{n!} + o(1)\).
    
    В достаточно малой окрестности точки \(x_0\) выражение \(\frac{f^{(n)}(x_0)}{n!} + o(1)\) сохраняет знак числа \(\frac{f^{(n)}(x_0)}{n!}\).
    \begin{itemize}
        \item Если \(f^{(n)}(x_0) > 0\), то \(f(x) - f(x_0) > 0\) для \(x \neq x_0\) \(\Rightarrow\) \(x_0\) — точка минимума.
        \item Если \(f^{(n)}(x_0) < 0\), то \(f(x) - f(x_0) < 0\) для \(x \neq x_0\) \(\Rightarrow\) \(x_0\) — точка максимума.
    \end{itemize}
    
    \textbf{Случай 2: \(n\) нечётное.}
    Тогда \((x - x_0)^n\) меняет знак при переходе через \(x_0\):
    \[
    (x - x_0)^n > 0 \text{ при } x > x_0, \quad (x - x_0)^n < 0 \text{ при } x < x_0.
    \]
    
    Следовательно, \(f(x) - f(x_0)\) также меняет знак, и \(x_0\) не является точкой экстремума (это точка перегиба).
\end{Proof}

\begin{consequence}
    На практике чаще всего используется частный случай при \(n = 2\):
    \begin{itemize}
        \item Если \(f'(x_0) = 0\) и \(f''(x_0) > 0\), то \(x_0\) — точка минимума.
        \item Если \(f'(x_0) = 0\) и \(f''(x_0) < 0\), то \(x_0\) — точка максимума.
    \end{itemize}
\end{consequence}

\begin{example}
    Исследуем функцию \(f(x) = x^4\) в точке \(x_0 = 0\):
    \begin{itemize}
        \item \(f'(0) = 0\), \(f''(0) = 0\), \(f'''(0) = 0\), \(f^{(4)}(0) = 24 > 0\).
        \item \(n = 4\) (чётное), \(f^{(4)}(0) > 0\) \(\Rightarrow\) \(x_0 = 0\) — точка минимума.
    \end{itemize}
    
    Для функции \(f(x) = x^3\):
    \begin{itemize}
        \item \(f'(0) = 0\), \(f''(0) = 0\), \(f'''(0) = 6 \neq 0\).
        \item \(n = 3\) (нечётное) \(\Rightarrow\) \(x_0 = 0\) — точка перегиба.
    \end{itemize}
\end{example}
